\subsection{GENERAL DEFINITION OF DERIVATIONS}

Consider a commutative ring K, six graded K-modules of type \(\Delta\) : A, A', A'', B, B', B'', and three K-linear mappings

\[
\mu \colon \mathrm{A} \times \mathrm{A} ^ {\prime} \rightarrow \mathrm{A} ^ {\prime \prime}, \quad \lambda_ {1} \colon \mathrm{B} \times \mathrm{A} ^ {\prime} \rightarrow \mathrm{B} ^ {\prime \prime}, \quad \lambda_ {2} \colon \mathrm{A} \times \mathrm{B} ^ {\prime} \rightarrow \mathrm{B} ^ {\prime \prime}
\]

such that the corresponding K-linear mappings

\[
\mathrm{A} \otimes_ {\mathrm{K}} \mathrm{A} ^ {\prime} \rightarrow \mathrm{A} ^ {\prime \prime}, \quad \mathrm{B} \otimes_ {\mathrm{K}} \mathrm{A} ^ {\prime} \rightarrow \mathrm{B} ^ {\prime \prime}, \quad \mathrm{A} \otimes_ {\mathrm{K}} \mathrm{B} ^ {\prime} \rightarrow \mathrm{B} ^ {\prime \prime}
\]

are graded of degree 0. The image \(\mu(a, a')\) for \(a \in A\) , \(a' \in A'\) is simply denoted by \(a, a'\) or even \(aa'\) , and similarly for the two other bilinear mappings. The degree of \(a, a'\) is therefore the sum of the degrees of \(a\) and \(a'\) .

DEFINITION 1. Given the above and a commutation factor \(\varepsilon\) on \(\Delta \times \Delta\) , an \(\varepsilon\) -derivation (or \((\mathbf{K}, \varepsilon)\) -derivation) of degree \(\delta \in \Delta\) of \((\mathbf{A}, \mathbf{A}', \mathbf{A}'')\) into \((\mathbf{B}, \mathbf{B}', \mathbf{B}'')\) is a triple of graded K-linear mappings of degree \(\delta\) :

\[
d \colon \mathrm{A} \rightarrow \mathrm{B}, \quad d ^ {\prime} \colon \mathrm{A} ^ {\prime} \rightarrow \mathrm{B} ^ {\prime}, \quad d ^ {\prime \prime} \colon \mathrm{A} ^ {\prime \prime} \rightarrow \mathrm{B} ^ {\prime \prime}
\]

such that, for every homogeneous element \(a \in A\) and every element \(a' \in A'\)

\[
d ^ {\prime \prime} (a. a ^ {\prime}) = (d a). a ^ {\prime} + \varepsilon_ {\delta , \deg (a)} a. (d ^ {\prime} a ^ {\prime}).\tag{4}
\]

It obviously suffices by linearity to verify relation (4) when \(a\) and \(a'\) run through respective generating systems of A and A'.

It is often convenient to denote the three mappings \(d, d', d''\) by the same letter \(d\) (which can be justified by denoting equally by \(d\) the graded K-linear mapping of degree \(\delta\) )

\[
(a, a ^ {\prime}, a ^ {\prime \prime}) \mapsto (d a, d ^ {\prime} a ^ {\prime}, d ^ {\prime \prime} a ^ {\prime \prime})
\]

of \(\mathbf{A} \oplus \mathbf{A}' \oplus \mathbf{A}''\) into \(\mathbf{B} \oplus \mathbf{B}' \oplus \mathbf{B}''\) . Relation (4) can then be written more simply

\[
d (a. a ^ {\prime}) = (d a). a ^ {\prime} + \varepsilon_ {\delta . \deg (a)} a. (d a ^ {\prime}).\tag{5}
\]

The \(\varepsilon\) -derivations of (A, A', A'') into (B, B', B'') of given degree form a sub-K-module of the K-module of graded linear mappings

\[
\operatorname{Homgr} _ {\mathbb {K}} (\mathrm{A} \oplus \mathrm{A} ^ {\prime} \oplus \mathrm{A} ^ {\prime \prime}, \mathrm{B} \oplus \mathrm{B} ^ {\prime} \oplus \mathrm{B} ^ {\prime \prime}).
\]

When \(\varepsilon(\alpha, \beta) = 1\) for all \(\alpha, \beta\) in \(\Delta\) , we say simply derivation (or K-derivation) instead of \(\varepsilon\) -derivation. Derivations form a sub-K-module of

\[
\operatorname{Hom} _ {\mathbb {K}} (\mathrm{A} \oplus \mathrm{A} ^ {\prime} \oplus \mathrm{A} ^ {\prime \prime}, \mathrm{B} \oplus \mathrm{B} ^ {\prime} \oplus \mathrm{B} ^ {\prime \prime}).
\]

When \(\Delta = Z\) and \(\varepsilon(p, q) = (-1)^{pq}\) , every \(\varepsilon\) -derivation of even degree is a derivation; every \(\varepsilon\) -derivation of odd degree is often called an antiderivation (or K-antiderivation); an antiderivation d therefore satisfies

\[
d (a. a ^ {\prime}) = (d a). a ^ {\prime} + (- 1) ^ {\deg (a)} a. (d a ^ {\prime})\tag{6}
\]

for a homogeneous element \(a \in A\) .

Remarks. (1) The notion of derivation can be defined for non-graded modules by agreeing to give these modules the trivial graduation.

\begin{enumerate}
\def\labelenumi{(\arabic{enumi})}
\setcounter{enumi}{1}
\tightlist
\item
  If only \(\varepsilon\) -derivations of given degree \(\delta\) are considered, the commutation factor \(\varepsilon\) may be disposed of as follows: the bilinear mapping \(\lambda_2: \mathbf{A} \times \mathbf{B}' \to \mathbf{B}''\) is modified by replacing it by the bilinear mapping \(\lambda_2': \mathbf{A} \times \mathbf{B}' \to \mathbf{B}''\) such that, for every homogeneous \(a\) in \(\mathbf{A}\) and all \(b' \in \mathbf{B}'\) ,
\end{enumerate}

\[
\lambda_ {2} ^ {\prime} (a, b ^ {\prime}) = \varepsilon_ {\delta , \deg (a)} \lambda_ {2} (a, b ^ {\prime}).
\]

Then \(d\) is a derivation relative to the bilinear mappings \(\mu, \lambda_{1}, \lambda_{2}^{\prime}\) .

The general definition of \(\varepsilon\) -derivations given above is especially used in two cases:

Case (I): A = B, \(A' = B'\) , \(A'' = B''\) and the three bilinear mappings \(\mu\) , \(\lambda_{1}\) , \(\lambda_{2}\) are equal to the same mapping.

Case (II): \(\mathbf{A} = \mathbf{A}' = \mathbf{A}''\) , \(\mathbf{B} = \mathbf{B}' = \mathbf{B}''\) , so that (for \(\mu\) ) \(\mathbf{A}\) is a graded algebra and the two K-bilinear mappings.

\[
\lambda_ {1} \colon \mathrm{B} \times \mathrm{A} \rightarrow \mathrm{B}, \quad \lambda_ {2} \colon \mathrm{A} \times \mathrm{B} \rightarrow \mathrm{B}\tag{7}
\]

are such that the corresponding K-linear mappings \(B \otimes_{K} A \to B\) , \(A \otimes_{K} B \to B\) are graded of degree 0. An \(\varepsilon\) -derivation of degree \(\delta\) of A into B is then a graded K-linear mapping \(d: A \to B\) of degree \(\delta\) , such that for every homogeneous x in A and every \(y \in A\) , we have the relation

\[
d (x y) = (d x) y + \varepsilon_ {\delta , \deg (a)} x (d y).\tag{8}
\]

Consider in particular in case (II) the case where A is a unital associative K-algebra and \(\lambda_{1}\) and \(\lambda_{2}\) are the external laws of an (A, A)-bimodule (§ 4, no. 3, Definition 3). This holds notably when A and B are two unital associative K-algebras, a unital homomorphism of graded K-algebras \(\rho: A \to B\) is given and an (A, A)-bimodule structure is considered on B defined by the two external laws

\[
\lambda_ {1} \colon (b, a) \mapsto b \rho (a), \quad \lambda_ {2} \colon (a, b) \mapsto \rho (a) b
\]

for \(a\in \mathbf{A},b\in \mathbf{B}.\)

Cases (I) and (II) have the following case in common: consider a graded K-algebra A, take B = A, mappings (7) both being multiplication on A. We then speak of an ε-derivation (or (K, ε)-derivation) of the graded algebra A: it is a graded K-linear mapping of A into itself, of degree δ, satisfying (8) for every homogeneous x in A and all \(y \in A\) . In particular if A is a graded ring, considered as an (associative) Z-algebra, we speak of the ε-derivation of the ring A.

Let A be a unital commutative associative K-algebra and B an A-module; when we speak of a derivation of A into B, it will always be understood that we mean with the A-bimodule structure on B derived from its A-module structure; then the formula

\[
d (x y) = x (d y) + y (d x) \quad \text { for } \quad x \in \mathrm{A}, y \in \mathrm{A}\tag{9}
\]

holds for such a derivation \(d: A \rightarrow B\) .

